\chapter*{Resumen}
Este proyecto de la asignatura Complementos de Bases de Datos desarrolla una solución
de analítica de datos masivos para movilidad urbana utilizando el conjunto de viajes de
taxi de Nueva York (\textit{NYC Taxi \& Limousine Commission}, Yellow Taxi).
La implementación integra Apache Spark como motor de procesamiento distribuido,
Apache Zeppelin como entorno de análisis interactivo y Python/PySpark para
enriquecimiento geográfico y visualización con mapas.

La aportación central es un flujo completo y reproducible: carga de ficheros Parquet,
limpieza trazable con reglas explícitas, enriquecimiento temporal y geográfico, y
ejecución de consultas analíticas en dos notebooks de Zeppelin. El primero,
\texttt{Análisis Taxis Nueva York.zpln}, concentra la carga, el preprocesamiento y
todos los análisis tabulares y gráficos. El segundo,
\texttt{Análisis Taxis Nueva York Mapas.zpln}, consume los indicadores exportados
por zona y genera mapas coropléticos interactivos con Folium.

Entre los hallazgos principales, el análisis muestra picos de actividad en franjas de
tarde-noche, un volumen claramente dominante de viajes estándar frente a viajes a
aeropuerto, y contrastes muy marcados entre distritos en velocidad media, tarifa y
equilibrio de flujos de movilidad.

Como limitaciones, el despliegue se realiza en infraestructura local y no en clúster
productivo. Aun así, la arquitectura deja preparada la migración a entornos más
escalables (Spark Standalone, YARN o Kubernetes) sin rehacer la lógica analítica.

\vspace{.5cm}

\textbf{Palabras clave:} Big Data, Apache Spark, Apache Zeppelin, PySpark, Parquet,
movilidad urbana, Spark SQL, analítica urbana, calidad del dato, Folium.